% Physical page 2 of the concise study, which carries this sheet and nothing
% else. The dates, relation labels, alternatives and unresolved states are the
% generated projection of research/chronology.toml, and stand exactly as the
% terminal appendix of the expansive study prints them; the Location cells and
% the explanatory rows are compressed for this page. The generated definitions
% are imported here, once, because the concise entrypoint's own preamble does
% not carry them.
\input{liturgy/roman-rite/postconciliar/roman-missal-third-edition-en-us-2011/propers/temporal/pc-s54-twenty-eighth-sunday-in-ordinary-time-year-a/research/chronology-annotations}

\section{Scriptural Date and Location}

\zlabel{triptych:concise:chronology:start}
\begin{concisedossiertable}
Responsorial Psalm & Ps 23 (22):1--3a, 3b--4, 5, 6; response v.~6cd & No place
named in the verses &
\chronodate{responsorial-psalm}{\chronologyannotation{responsorial-psalm}}\\
\cmidrule(lr){1-4}
\dossierprose{Inherited attribution: the Vulgate title, in the Douay--Rheims
``A psalm for David'' (22:1), ascribes the psalm to David and gives no date or
occasion, so the setting answer is unresolved. The composition date is the
Psalter's common critical bound, from the introduction to the Psalms in the
New American Bible Revised Edition, which holds that no psalm can be dated
securely; it is not a date for this psalm.}
\midrule
Communion Antiphon, first option & cf.\ Ps 34 (33):11, an adaptation & By the
title, the court of Geth, where David fled from Saul to Achis (1~Kgs 21:10) &
\chronodate{communion-antiphon-a}{\chronologyannotation{communion-antiphon-a}}\\
\cmidrule(lr){1-4}
\dossierevent{Narrated event, by the title: David changing his countenance
before Achimelech (33:1), referred by the Douay--Rheims to 1~Kings 21:10--15.}
\cmidrule(lr){1-4}
\dossierprose{Inherited attribution: the title ``For David''. The attribution
answer is David's reign in the usual chronology the \work{Catholic
Encyclopedia} reports (Corbett), marked disputed. The superscription answer
dates the episode the title names, on Haydock's A.M. reckoning; it is not a
date of composition. The composition answer is the same common bound. The
antiphon's ``rich'' is the Latin and Greek reading; the Hebrew has young
lions.}
\midrule
Entrance Antiphon & Ps 130 (129):3--4, an adaptation & No place named in the
verses &
\chronodate{entrance-antiphon}{\chronologyannotation{entrance-antiphon}}\\
\cmidrule(lr){1-4}
\dossierprose{A song of ascents, ``A gradual canticle''; the title names no
author and no occasion. The composition answer is the common bound, not a date
for this psalm.}
\midrule
First Reading & Is 25:6--10a & ``In this mountain'' (25:6, 10), the ``mount
Sion'' and Jerusalem where the Lord of hosts reigns (24:23) &
\chronodate{first-reading}{\chronologyannotation{first-reading}}\\
\cmidrule(lr){1-4}
\dossierprose{Traditional attribution: ``Isaias the son of Amos'' (1:1). The
answer is the length of the prophet's ministry in the \work{Catholic
Encyclopedia}'s article on Isaias (Souvay); it dates the prophet, not
chapters 24--27 or this oracle, for which no composition date is held.}
\midrule
Gospel, longer form & Mt 22:1--14 & In the temple at Jerusalem, to ``the chief
priests and ancients of the people'' (21:23) &
\chronodate{gospel-long}{\chronologyannotation{gospel-long}}\\
Gospel, shorter form & Mt 22:1--10 & As for the longer form &
\chronodate{gospel-short}{\chronologyannotation{gospel-short}}\\
\cmidrule(lr){1-4}
\dossierevent{Narrated event: a parable spoken in the temple, after the
question about Jesus' authority and the parables of the two sons and the
vineyard.}
\cmidrule(lr){1-4}
\dossierprose{Traditional attribution: St~Matthew the Apostle. The event answer
dates the telling of the parable, not the wedding it narrates: a derivation
within Maas's chronology in the \work{Catholic Encyclopedia}, which places it
on the Tuesday of the last week in Jerusalem. The composition figures are
those the \work{Catholic Encyclopedia} reports: Jacquier's, from ancient
writers ``at variance'' and from Catholic critics where ``opinion is rather
divided'', and last Durand's, for an original text he holds to have been
Aramaic. They date the book, not the saying.}
\midrule
Alleluia verse & cf.\ Eph 1:17--18, an adaptation & From Paul, ``the prisoner
of Jesus Christ'' (3:1); to the saints ``at Ephesus'' (1:1) &
\chronodate{gospel-acclamation}{\chronologyannotation{gospel-acclamation}}\\
\cmidrule(lr){1-4}
\dossierprose{Traditional attribution: St~Paul, writing in captivity. The range
is the \work{Catholic Encyclopedia}'s for the letters of the captivity, whose
place, Caesarea or Rome, it calls ``a much mooted question''; the year is
Prat's. They date the letter, not the verse made from it.}
\midrule
Second Reading & Phil 4:12--14, 19--20 & From Paul ``in my bands'' (1:7, 13),
with greetings from ``Caesar's household'' (4:22); to the saints ``at
Philippi'' (1:1) &
\chronodate{second-reading}{\chronologyannotation{second-reading}}\\
\cmidrule(lr){1-4}
\dossierprose{Traditional attribution: St~Paul, writing in captivity. The year
is Prat's for the letters of the captivity; the range is Vander Heeren's for
Paul ``at Rome'', which he calls ``the nearly unanimous opinion''.}
\midrule
Communion Antiphon, second option & 1~Jn 3:2, an adaptation & No place named
in the letter &
\chronodate{communion-antiphon-b}{\chronologyannotation{communion-antiphon-b}}\\
\cmidrule(lr){1-4}
\dossierprose{Traditional attribution: St~John the Apostle. The range is
Durand's for ``the Johannine writings''; Drum finds the arguments ``probable
in favour of Ephesus and also for the last few years of the first century''.}
\end{concisedossiertable}
\zlabel{triptych:concise:chronology:end}
